\subsection{The exact global Jacobian class selected at two}

\begin{statement}{A.190}{Proved}
{Injective completed localization and the four-point two-adic Selmer intersection}
Let \(J_8=\operatorname{Jac}(C_8)\),
\(C_8:3u^8-2v^8=1\), and let
\(j:C_8\to J_8\) be \(R\mapsto[R-P_0]\),
where \(P_0=(1,1)\).  The actual localization map
induced by \(J_8(\Q)\hookrightarrow J_8(\Q_2)\)
is injective:
\[
 \widehat{J_8(\Q)}_2
         \longrightarrow\widehat{J_8(\Q_2)}_2.
\tag{A.190.1}
\]
The same is true for each of the twelve factors in
the particular decomposition (A.182.3).
With the continuous Selmer group of (A.187.5), the
exact intersection, including its global coordinate,
is
\[
\begin{aligned}
\bigl\{(\eta,R)\in&\operatorname{Sel}(\Q,T_2J_8)
                         \times C_8(\Q_2):\\
       &\operatorname{loc}_2\eta=\kappa_2(j(R))\bigr\}
   =\bigl\{(\kappa_2(j(P)),P):
                   P\in\{(\pm1,\pm1)\}\bigr\}.
\end{aligned}
\tag{A.190.2}
\]
The four global classes on the right are distinct.

In particular, for the full metabelian sections in
(A.189.5), the fixed-sheet map to the compact
two-adic Jacobian extension has global class exactly
\(\kappa_2(j(P))\), where \(P\) is that sheet's
dyadic point in \(\{(\pm1,\pm1)\}\).
The map and the equality are in the stated continuous
two-adic compact-Jacobian quotient.  The full
classification remains the open assertion A1.
\end{statement}

\begin{proof}
\emph{Local completions and their finite torsion.}
For an abelian variety over a finite extension of
\(\Q_2\), its local group is compact.  Its abelian
logarithm is a homomorphism which is an analytic
isomorphism on a sufficiently small open subgroup,
and its image there is a full \(\Z_2\)-lattice in
the Lie algebra.  These are the local-group and
formal-antiderivative comparisons in
\cite[\S4.1 and Remark 4.3, author version pp.3--5]{LegMcCallumPoonenChabauty}.
The same proof works over every finite extension:
the formal coordinates and convergent primitives
are taken in its valuation ring.  Its logarithm
kernel is discrete and compact, hence finite;
it is exactly its torsion subgroup.

The natural map from the local group to its
two-primary completion has kernel precisely its
odd-order torsion.  Here is a direct verification
of the kernel.  An element in \(\bigcap_n2^nA(k)\)
has logarithm in every \(2^n\log A(k)\); the
bounded \(\Z_2\)-lattice \(\log A(k)\) has
zero intersection of those subgroups.  The element
is consequently torsion.  All its preimages under
\(2^n\) are torsion too, since their logarithms
are zero.  In the finite torsion group the
intersection is its odd-order subgroup,
since multiplication by two is invertible exactly
on that subgroup.  Conversely that subgroup lies
in every \(2^nA(k)\).  Compactness also makes the
map onto the inverse limit surjective: a compatible
system of cosets has a nonempty intersection.
This proves the asserted description of the
completion.  In particular ordinary rational
two-primary torsion injects into the local
completion.

For every rank-one factor in A.185--A.186, the
displayed rational point of infinite order remains
of infinite order in its local group: a torsion
equation there would be the same rational point
equation.  Its logarithm is therefore nonzero.
The continuous map from its completed free
\(\Z_2\)-part is injective after applying logarithm;
the remaining rational two-primary torsion is
injective by the preceding paragraph.  This proves
completed localization injective for the three
rational elliptics, all six genus-two factors,
and the rank-one elliptic over \(\Q(\sqrt{-3})\).
For the latter field, two has a unique place:
the polynomial of a primitive cube root is
\(X^2+X+1\), irreducible modulo two.

\emph{The rank-two elliptic over \(\Q(\sqrt2)\).}
Put \(s^2=2\), \(k=\Q(s)\),
\(E:y^2=x^3+6sx\), and
\(P=(3,3+3s)\).  Let \(\sigma(s)=-s\).
The coefficient curve of \(\sigma P\) is
\(E^\sigma:y^2=x^3-6sx\).  Use the actual
map
\[
 \beta:E^\sigma\longrightarrow E,
 \qquad
 \beta=\psi_{6s}\circ\lambda_s,
 \quad\lambda_s(x,y)=(s^2x,s^3y),
 \qquad Q=\beta(\sigma P).
\tag{A.190.3}
\]
The scaling has target \(E_{-24s}\), the domain
of \(\psi_{6s}\).  For the chosen invariant
forms \(\omega=dx/(2y)\), direct differentiation
of (A.185.8) gives
\(\psi_{6s}^*\omega_E=2\omega_{E_{-24s}}\).
The scaling then gives
\(\beta^*\omega_E=(2/s)\omega_{E^\sigma}
 =s\omega_{E^\sigma}\).

There is a unique dyadic place, with completion
\(k_2=\Q_2(s)\), since \(X^2-2\) is
Eisenstein.  Write the scalar elliptic logarithm
of \(P\), evaluated on the specified form, as
\(\ell(P)=a+bs\), with \(a,b\in\Q_2\).
Naturality under the displayed map and under
\(\sigma\) gives
\[
                  \ell(Q)=s\sigma(\ell(P))=-2b+as.
\tag{A.190.4}
\]
The point \(P\) is of infinite order by A.182--A.185,
so \((a,b)\ne(0,0)\).  The determinant of these
two coefficient vectors is \(a^2+2b^2\).
When both terms are nonzero, their two-adic
valuations have different parity, and they cannot
sum to zero; when one is zero the assertion is
immediate.  Thus the two local logarithms are
\(\Q_2\)-independent.  The global rank is two
by A.185, so the logarithm of the full global
free part has rank two.  With the rational torsion
argument above, this proves
\(\widehat{E(k)}_2\to\widehat{E(k_2)}_2\)
injective.

\emph{The rank-two elliptic over \(\Q(\sqrt{-6})\).}
Now put \(s^2=-6\), \(k=\Q(s)\), and
\(E:y^2=x^3+8sx\).  Two actual global points
used in the computation are
\[
\begin{aligned}
 P&=(-2,-4+2s),\qquad R=(-5+4s,37-s)\in E_{-2s}(k),\\
 Q&=\psi_{8s}(4x(R),8y(R))\\
  &=\bigl(x(R)-2s/x(R),
              y(R)(1+2s/x(R)^2)\bigr)\in E(k).
\end{aligned}
\tag{A.190.5}
\]
The first two point equations were proved in
A.185.  The last line is the literal dual
isogeny after its indicated scaling.  The dyadic
completion is \(k_2=\Q_2(s)\); the polynomial
\(X^2+6\) is Eisenstein.  Use the extension
of \(v_2\) with \(v_2(s)=1/2\).

We prove nonvanishing of its logarithm determinant
with an explicit error bound.  For a point on
\(E_b:y^2=x^3+bx\) near infinity, put
\(t=-x/y\).  The root with a double pole is
\[
 x(t)=\frac{1+\sqrt{1-4bt^4}}{2t^2},
 \qquad y(t)=-x(t)/t,
 \qquad \frac{dx}{2y}=\frac{dt}{\sqrt{1-4bt^4}}.
\tag{A.190.6}
\]
Both identities follow from
\(x^2-x/t^2+b=0\) and differentiation.
Termwise integration gives the formal logarithm
for this exact invariant differential:
\[
 \ell_b(t)=\sum_{n\geq0}
           \binom{2n}{n}\frac{b^nt^{4n+1}}{4n+1}.
\tag{A.190.7}
\]
Its constant is zero.  Translation invariance
and the formal chain rule identify it with
the abelian logarithm on this open subgroup.
All denominators \(4n+1\) are odd.

Use the exact duplication formulas
\[
 m=\frac{3x^2+b}{2y},\qquad
 x(2Z)=m^2-2x,\qquad
 y(2Z)=m(x-x(2Z))-y.
\tag{A.190.8}
\]
For \(b=8s\), two applications to each point
in (A.190.5) give
\[
\begin{array}{c|c|c|c}
Z&v_2(t(4Z))&v_2(x(4Z))&v_2(y(4Z))\\ \hline
P&3/2&-3&-9/2\\
Q&2&-4&-6.
\end{array}
\tag{A.190.9}
\]
Thus both points lie on the pole branch in
(A.190.6).  For each one put \(t_Z=t(4Z)\)
and
\(L_Z=(t_Z+2bt_Z^5/5)/4\).
The omitted terms in \(\ell(Z)-L_Z\) have
valuation at least
\[
 2v_2(b)+9v_2(t_Z)-2
        =\begin{cases}37/2,&Z=P,\\23,&Z=Q.\end{cases}
\tag{A.190.10}
\]
Indeed their indices satisfy \(n\geq2\),
the central binomial coefficient is integral,
and \(v_2(b)=7/2\); the displayed lower
bound increases with \(n\).

For a field element \(a+bs\), its valuation
is \(\min(v_2(a),v_2(b)+1/2)\); the two
possible finite values have different parity
after multiplication by two.  The error bounds
in (A.190.10) therefore bound the two coefficient
errors for \(P\) by \((19,18)\), and those
for \(Q\) by \((23,23)\).  Exact rational
calculation of the two retained terms gives
\[
 \begin{pmatrix}
  [1](2\ell(P))&[1](2\ell(Q))\\
  [s](2\ell(P))&[s](2\ell(Q))
 \end{pmatrix}
 \equiv
 \begin{pmatrix}2526&270\\1435&4006\end{pmatrix}
                  \pmod{4096}.
\tag{A.190.11}
\]
Here \([1]\) and \([s]\) mean the literal
coefficient maps in the basis \((1,s)\), and
the factor two is multiplication of the logarithm
by two.  The stated tail bounds justify every
entry of this congruence.  Its determinant is
two modulo four, so the determinant of the
unmultiplied logarithms has valuation \(-1\)
and is nonzero.  The exact certificate also
checks that the determinant-error valuation is
at least eighteen.  Hence these two logarithms
are \(\Q_2\)-independent.  The rank-two result
of A.185 and the rational torsion argument prove
completed localization injective for this elliptic
curve as well.

\emph{Transfer through the actual rational isogenies.}
For the three Weil restrictions in A.185,
the equalities of groups of rational and local
points are respectively \(E(k)\) and \(E(k_2)\).
They induce the same equalities on their
two-primary completions.  All three completed
localization maps are therefore injective by
the preceding calculations.

Let \(f_{\mathcal P},g_{\mathcal P}\) be the
actual rational Prym maps in (A.187.12)--(A.187.13),
with \(g_{\mathcal P}f_{\mathcal P}=[2]\).
A completed Prym class with zero localization
maps to zero in the global Weil-restriction
completion by its injectivity.  Multiplication
by two kills the original class.  Its class
is consequently ordinary rational two-primary
torsion, since the global group is finitely
generated.  The already proved torsion
injection makes it zero.  This proves the
three Prym assertions.

The product \(\mathcal B\) has injective
completed localization.  A class \(z\) in
the kernel of (A.190.1) therefore has \(Fz=0\).
Applying the literal reverse map of A.188
gives \(16z=0\).  This is again rational
two-primary torsion, whose localization is
injective.  Thus \(z=0\), proving (A.190.1)
with the asserted maps.

\emph{The four global Selmer classes and the section map.}
By A.187 every \(\eta\) on the left of
(A.190.2) has a unique global completed
Mordell--Weil preimage.  Push it forward
along the first genus-two map and add
\(D=[(3,9)-\infty]\).  The equality of
local Kummer classes then becomes the
condition of (A.189.2) for the image of
\(R\) under (A.189.3).  Thus A.189 makes
\(R\) one of the four rational points.
Its rational class \(j(R)\) has the same
localization as the global preimage of
\(\eta\).  Injectivity of (A.190.1)
identifies the two preimages and proves
the asserted equality in continuous
cohomology.  Every listed pair satisfies
that equality by functoriality, proving
both directions of (A.190.2).

The four classes are separated by the second
and third genus-two components of \(F\).
At \(P=(\epsilon_u,\epsilon_v)\), their
images, with the base point \(P_0\)
subtracted, are
\[
       (\epsilon_v-1)D_{648},\qquad
       (\epsilon_u-1)D_{-432}.
\tag{A.190.12}
\]
This follows directly from their coordinates
\((6v^2,108u^4v)\) and
\((6u^2,72uv^4)\).  The two displayed
classes \(D_c\) generate free \(\Z_2\)
summands by A.186.  The pair of images
therefore separates all four sign choices.

For the metabelian statement, use exactly
the fixed-sheet cover extension (A.189.18).
Compactification and \(j:C_8\to J_8\)
give its equivariant geometric map
\(H^{\mathrm{ab}}\to T_2J_8\) and its
arithmetic extension map.  The ordinary
and cuspidal fixed-sheet comparisons in
A.189 put each finite image in the
corresponding global Selmer group; the
real-place component is included there.
The continuous comparison of A.187
gives a class \(\eta\) as in (A.190.2),
with precisely that sheet's dyadic point.
Applying (A.190.2) proves the stated
global compact coefficient equality.

All rational quadratic-field arithmetic,
the actual dual-isogeny point, both exact
duplications, the formal-branch valuations,
and the determinant with its error bounds
are retained in
\path{code/point3_fermat8_dyadic_regulator_certificate.py}
and its output and rational data.  The
infinite-series bound is (A.190.10).
\end{proof}
